%!TEX root = ../graduate-work.tex
\section{Verification Algorithms and Tools}

\subsection{Hierarchy of Bound-Propagation Methods}

The foundation of modern verifiers is bound propagation---algorithms that compute,
for each network output, an interval guaranteed to contain the true value for any
admissible input. Three levels of accuracy and computational cost are distinguished.

\textbf{IBP (Interval Bound Propagation)} propagates intervals $[\mathbf{l}, \mathbf{u}]$
layer by layer through the network.
For a linear layer $\mathbf{y} = W\mathbf{x} + \mathbf{b}$:
$$\mathbf{l}_y = W^+\mathbf{l}_x + W^-\mathbf{u}_x + \mathbf{b}, \quad
  \mathbf{u}_y = W^+\mathbf{u}_x + W^-\mathbf{l}_x + \mathbf{b},$$
where $W^+ = \max(W, 0)$ and $W^- = \min(W, 0)$ elementwise.
The method is cheap ($O(n)$, a single forward pass), but intervals widen
rapidly with depth due to accumulating over-approximation
error~\cite{CROWNPaper}.

\textbf{CROWN (Convex Relaxation with Optimised Neurons)} computes bounds via
a backward pass analogous to backpropagation.
For a ReLU neuron with pre-activation $x \in [l, u]$, $l < 0 < u$, a linear
relaxation is constructed:
$$\lambda_l x + b_l \leq \mathrm{ReLU}(x) \leq \frac{u}{u - l}(x - l),$$
where $\lambda_l \in [0, 1]$ is fixed heuristically.
The backward pass links the network output to the input through a chain of
linear inequalities, yielding substantially tighter bounds than
IBP~\cite{CROWNPaper}.

\textbf{$\alpha$-CROWN} takes this further: the slopes $\lambda_l$ become
per-neuron optimisable parameters $\alpha_i$. Gradient descent over $\alpha$
raises the lower bound (or lowers the upper bound), producing the tightest
possible bounds without any branching~\cite{AlphaBetaCrownVNNCOMP23}.

\subsubsection{Incomplete and Complete Verification}

IBP, CROWN, and $\alpha$-CROWN are methods for \textit{incomplete verification}.
They construct an over-approximation of the output set, satisfying
\textit{soundness}: lower bounds never exceed the true minima, and upper bounds
are never below the true maxima. When the resulting bounds are tight enough to
confirm the property, verification succeeds; otherwise the outcome is
indeterminate. Crucially, no false positive can occur: if the property is
violated for even one input, the bounds will reflect it.

\textit{Complete verification} guarantees a definitive answer: either the
property holds, or a counterexample is produced. To achieve this,
\textit{Branch-and-Bound} (BaB) is applied: when bounds are insufficiently
tight, the input space is split into sub-regions (e.g.\ by fixing a ReLU state:
$x \geq 0$ or $x < 0$) and bounds are recomputed for each sub-region. The
process repeats until an answer is obtained or resources are exhausted.

\subsection{Alpha-Beta-CROWN}

$\alpha,\beta$-CROWN~\cite{AlphaBetaCrownVNNCOMP23}---winner of VNN-COMP
from 2021 through 2024---combines the above methods into a unified verifier.
For incomplete verification it uses $\alpha$-CROWN; when that does not yield
sufficiently tight bounds it falls back to complete verification via BaB.

\textbf{$\beta$-CROWN: complete verification without LP.}
Traditional BaB-based verifiers (e.g.\ MIPVerify) solve a linear program at
each leaf of the search tree---this is slow and requires data transfer from GPU
to CPU. $\beta$-CROWN encodes branching constraints (``neuron $i > 0$'' or
``neuron $i \leq 0$'') directly into the bound-propagation procedure via
Lagrangian dual variables ($\beta$). As a result, bound evaluation for
thousands of branches is performed in parallel on the GPU as a series of
matrix multiplications, and $\beta$ parameters are optimised jointly with
$\alpha$ via gradient descent.

\subsubsection{Cutting-Plane Methods}

To strengthen relaxations without full branching, cutting planes from integer
programming are employed. GCP-CROWN adds linear inequalities that jointly
constrain multiple neurons, cutting off fractional solutions that are feasible
in the LP relaxation but infeasible for the actual ReLU network. BICCOS
(Branch-and-bound Inferred Cuts with COnstraint Strengthening) generates cuts
from logical implications in the search tree and applies them iteratively to
tighten bounds at root nodes---enabling problems to be solved that would
otherwise require prohibitively deep
branching~\cite{AutoLiRPAPracticeSlides}.

\subsection{Other Verifiers}

\textbf{VeriNet} uses Symbolic Interval Propagation (SIP): instead of linear
inequalities, CROWN-style dependencies are maintained as symbolic expressions.
The tool is particularly effective for input splitting---a strategy that works
well for low-dimensional inputs, as in control tasks (ACAS~Xu)~\cite{IntervalReachabilityAnalysis}.

\textbf{nnenum} (winner of VNN-COMP 2020) dynamically switches between levels
of set representation precision---from fast but coarse zonotopes to accurate
but expensive ImageStars---depending on the difficulty of the current
sub-problem~\cite{NNEnumPaper}.

\textbf{Marabou~2.0} addresses the same problem differently: instead of binary
search over neuron states, it employs the DeepSoI procedure
(Sum of Infeasibilities)---a continuous function measuring the total violation
of ReLU constraints, which can be minimised by gradient
methods. This transforms the combinatorial search into an optimisation
problem~\cite{arXiv2401Marabou2}.

\textbf{NeuralSAT} adapts CDCL (Conflict-Driven Clause Learning) from SAT
solving: upon detecting an infeasible combination of neuron states, it records
it as a ``nogood'', and entire subtrees are pruned without revisitation in
subsequent search~\cite{NeuralSATRepo}.

\subsection{The \texorpdfstring{\texttt{auto\_LiRPA}}{auto\_LiRPA} Library}

All experiments in this work are based on
\texttt{auto\_LiRPA}~\cite{AlphaBetaCrownVNNCOMP23}---the library on which
$\alpha,\beta$-CROWN is built. Its operating principle: bound-propagation
operations are described as nodes of a PyTorch computational graph, and the
library automatically derives the graph for upper and lower bounds of any
intermediate network node. This enables verification of arbitrary architectures
(CNN, ResNet, Transformer) without writing specialised
code~\cite{LiRPAVerifyRepo}.

A critically important feature is \textit{batch BaB}: instead of processing
branches sequentially, the verifier batches hundreds of thousands of unsolved
sub-problems and computes bounds for all of them in a single GPU pass.
This is precisely what makes implementing parallelism in \texttt{auto\_LiRPA}
non-trivial: any modification must preserve the ability to process batches
and remain compatible with the PyTorch computational graph.
